% =========================================================================
% SECCIÓN 8: ANEXOS (CARTOGRAFÍA HÍBRIDA MULTI-ESCENARIO)
% =========================================================================
\newpage
\section{ANEXO GRÁFICO: ESCENARIOS DE INVERSIÓN Y CARTOGRAFÍA HÍBRIDA}

La representación cartográfica de las matrices de deformación demanda una estratificación visual rigurosa según el anclaje espacial utilizado durante la inversión matemática de la red SBAS. A continuación, se despliega el catálogo completo de \textit{layouts} híbridos generados automatizadamente. En estos documentos planimétricos, la capa de InSAR se superpone sobre la ortofotografía base en escala de grises para resaltar el contraste estructural, e incluye un recuadro de contexto urbano. Se presentan los tres escenarios geodésicos evaluados para garantizar la transparencia y reproducibilidad de los resultados.

% -------------------------------------------------------------------------
\subsection{Escenario de Referencia A: Castillo de San Felipe}

En esta iteración algorítmica, la red asume como punto de velocidad nula ($0.0 \text{ mm/año}$) el polígono de cimentación correspondiente al Castillo de San Felipe de Barajas, ubicado estratégicamente al oriente del Centro Histórico sobre el cerro de San Lázaro. Resulta de excepcional utilidad para evaluar cómo se asienta la mancha urbana circundante frente al monumento patrimonial estructuralmente más estable de la ciudad.

\begin{figure}[H]
    \centering
    \includegraphics[width=1.0\linewidth]{LAYOUT_CASTILLO_Centro_Historico.jpg}
    \caption{Cartografía del Centro Histórico (Ref: Castillo). Ubicado en el núcleo fundacional de la ciudad (Patrimonio de la Humanidad), el panel lateral muestra la tendencia de deformación en el periodo analizado frente a la cimentación de la fortaleza.}
    \label{fig:castillo_historico}
\end{figure}

\begin{figure}[H]
    \centering
    \includegraphics[width=1.0\linewidth]{LAYOUT_CASTILLO_San_Diego.jpg}
    \caption{Detalle interferométrico del Barrio San Diego (Ref: Castillo). Este sector nororiental del recinto amurallado, de uso residencial histórico, permite observar diferencias espaciales en la velocidad de deformación periférica.}
    \label{fig:castillo_sandiego}
\end{figure}

\begin{figure}[H]
    \centering
    \includegraphics[width=1.0\linewidth]{LAYOUT_CASTILLO_Bocagrande.jpg}
    \caption{Despliegue geoespacial de Bocagrande (Ref: Castillo). Esta península costera está altamente densificada con infraestructura en altura erigida sobre antiguos depósitos litorales. Los vectores negativos advierten sobre la respuesta a la sobrecarga.}
    \label{fig:castillo_bocagrande}
\end{figure}

\begin{figure}[H]
    \centering
    \includegraphics[width=1.0\linewidth]{LAYOUT_CASTILLO_Barrio_Manga.jpg}
    \caption{Mapa cinemático del Barrio Manga (Ref: Castillo). Formación insular residencial conectada por puentes al núcleo histórico. El patrón espacial observado es consistente con la compresibilidad de los sustratos arcillosos blandos ante la redensificación reciente.}
    \label{fig:castillo_manga}
\end{figure}

\begin{figure}[H]
    \centering
    \includegraphics[width=1.0\linewidth]{LAYOUT_CASTILLO_GNSS_CART.jpg}
    \caption{Validación puntual de la antena GNSS CART00COL (Ref: Castillo). Ubicada en la Isla Manzanillo (CIOH), el encuadre permite evaluar el comportamiento de la losa de instalación frente a la estabilidad teórica asumida en la fortaleza militar.}
    \label{fig:castillo_gnss}
\end{figure}

\begin{figure}[H]
    \centering
    \includegraphics[width=1.0\linewidth]{LAYOUT_CASTILLO_Max_Subsidencia.jpg}
    \caption{Mapa del sector con mayor subsidencia estimada (Ref: Castillo). Este encuadre focaliza la coordenada geométrica sometida al mayor esfuerzo cortante de hundimiento transversal a toda el área de estudio.}
    \label{fig:castillo_maxsub}
\end{figure}

\begin{figure}[H]
    \centering
    \includegraphics[width=1.0\linewidth]{LAYOUT_CASTILLO_Max_Elevacion.jpg}
    \caption{Mapa del sector con máxima anomalía positiva (Ref: Castillo). Identificación espacial de las zonas que registran dinámicas contrarias a la tendencia general por el levantamiento local de materiales.}
    \label{fig:castillo_maxele}
\end{figure}


% -------------------------------------------------------------------------
\subsection{Escenario de Referencia B: Zona Elevada (Cerro)}

Este modelo de inversión asume estabilidad topológica en un vértice emplazado en la zona oriental montañosa (Cerro de La Popa), con una coherencia espacial sostenida superior a 0.9. Al separar la referencia geométrica del sustrato sedimentario costero, proporciona un marco de comparación de alto contraste frente a la planicie urbana.

\begin{figure}[H]
    \centering
    \includegraphics[width=1.0\linewidth]{LAYOUT_CERRO_Centro_Historico.jpg}
    \caption{Evaluación del polígono histórico colonial (Ref: Cerro). El evidente contraste de deformación subraya el asentamiento continuo de los depósitos marinos del nivel del mar frente al afloramiento rocoso oriental.}
    \label{fig:cerro_historico}
\end{figure}

\begin{figure}[H]
    \centering
    \includegraphics[width=1.0\linewidth]{LAYOUT_CERRO_San_Diego.jpg}
    \caption{Comportamiento cinemático en San Diego (Ref: Cerro). Las series temporales trazadas confirman que la franja de mampostería colonial más expuesta al mar Caribe experimenta dinámicas de hundimiento progresivo.}
    \label{fig:cerro_sandiego}
\end{figure}

\begin{figure}[H]
    \centering
    \includegraphics[width=1.0\linewidth]{LAYOUT_CERRO_Bocagrande.jpg}
    \caption{Monitoreo de deformación en Bocagrande (Ref: Cerro). El gradiente térmico del modelo expone la severidad de la compresión edáfica a lo largo del eje longitudinal de esta formación arenosa densificada.}
    \label{fig:cerro_bocagrande}
\end{figure}

\begin{figure}[H]
    \centering
    \includegraphics[width=1.0\linewidth]{LAYOUT_CERRO_Barrio_Manga.jpg}
    \caption{Distribución de hundimientos en Manga (Ref: Cerro). El análisis visual ratifica la vulnerabilidad estructural en el sur-oriente de esta isla frente a variaciones de la presión de poros en suelos finos.}
    \label{fig:cerro_manga}
\end{figure}

\begin{figure}[H]
    \centering
    \includegraphics[width=1.0\linewidth]{LAYOUT_CERRO_GNSS_CART.jpg}
    \caption{Métrica de discrepancia geodésica en estación CART00COL (Isla Manzanillo, Ref: Cerro). Extrae de forma puntual la diferencia matemática entre la red relativa InSAR anclada a la cota alta y la medición satelital.}
    \label{fig:cerro_gnss}
\end{figure}

\begin{figure}[H]
    \centering
    \includegraphics[width=1.0\linewidth]{LAYOUT_CERRO_Max_Subsidencia.jpg}
    \caption{Identificación espacial de la subsidencia crítica (Ref: Cerro). Cartografía del píxel en la periferia de la bahía con la más alta tasa de declive altimétrico continuo reportada por el sensor europeo.}
    \label{fig:cerro_maxsub}
\end{figure}

\begin{figure}[H]
    \centering
    \includegraphics[width=1.0\linewidth]{LAYOUT_CERRO_Max_Elevacion.jpg}
    \caption{Anomalía positiva de deformación vertical (Ref: Cerro). Localización topológica de dispersores aislados que exhiben vectores ascendentes dentro de la malla interferométrica analizada.}
    \label{fig:cerro_maxele}
\end{figure}


% -------------------------------------------------------------------------
\subsection{Escenario de Referencia C: Literatura Previa (Paper)}

Esta tercera calibración adopta las coordenadas de anclaje validadas exhaustivamente en literatura geodésica regional sobre el sector de la Isla Manzanillo (CIOH). Permite blindar la comparabilidad cronológica del estudio y descarta sesgos de marco que puedan introducirse por la geometría contemporánea.

\begin{figure}[H]
    \centering
    \includegraphics[width=1.0\linewidth]{LAYOUT_PAPER_Centro_Historico.jpg}
    \caption{Validación cruzada sobre el recinto amurallado fundacional (Ref: Literatura). La consistencia de los patrones de asentamiento respecto a la métrica externa consolida la confiabilidad geofísica del análisis.}
    \label{fig:paper_historico}
\end{figure}

\begin{figure}[H]
    \centering
    \includegraphics[width=1.0\linewidth]{LAYOUT_PAPER_San_Diego.jpg}
    \caption{Mapeo de vulnerabilidad predial en el nororiente amurallado de San Diego (Ref: Literatura). Provee una capa de priorización objetiva para canalizar futuras obras de apuntalamiento patrimonial.}
    \label{fig:paper_sandiego}
\end{figure}

\begin{figure}[H]
    \centering
    \includegraphics[width=1.0\linewidth]{LAYOUT_PAPER_Bocagrande.jpg}
    \caption{Firma interferométrica en la península de Bocagrande (Ref: Literatura). Persistencia del patrón de descenso, indicando un comportamiento no elástico e irreversible de los rellenos litorales.}
    \label{fig:paper_bocagrande}
\end{figure}

\begin{figure}[H]
    \centering
    \includegraphics[width=1.0\linewidth]{LAYOUT_PAPER_Barrio_Manga.jpg}
    \caption{Respuesta cinemática del polo residencial-portuario del Barrio Manga (Ref: Literatura). Sugiere una relación intrínseca entre la densificación urbana acelerada y las tasas de subsidencia marginal observadas.}
    \label{fig:paper_manga}
\end{figure}

\begin{figure}[H]
    \centering
    \includegraphics[width=1.0\linewidth]{LAYOUT_PAPER_GNSS_CART.jpg}
    \caption{Anclaje a estación SIRGAS CART00COL en Isla Manzanillo (Ref: Literatura). Extracción del vector de deformación específico en la losa de montaje, insumo de contraste vital para corroborar el sesgo de referencia.}
    \label{fig:paper_gnss}
\end{figure}

\begin{figure}[H]
    \centering
    \includegraphics[width=1.0\linewidth]{LAYOUT_PAPER_Max_Subsidencia.jpg}
    \caption{Vectores límite de subsidencia urbana (Ref: Literatura). Extracción del rango inferior de la serie de datos para documentar el escenario cinemático y estructural más conservador en la ciudad.}
    \label{fig:paper_maxsub}
\end{figure}

\begin{figure}[H]
    \centering
    \includegraphics[width=1.0\linewidth]{LAYOUT_PAPER_Max_Elevacion.jpg}
    \caption{Vectores límite de levantamiento superficial (Ref: Literatura). Rastreo espacial de los reflectores de radar cuyo comportamiento anómalo señala rebote elástico puntual.}
    \label{fig:paper_maxele}
\end{figure}
